\chapter{Algebra of curvature}

The curvature of a Riemannian manifold is described by a tensor, not just a number.
This is one of the principle differences between differential geometry of surfaces and higher-dimensional differential geometry.

In this lecture we will give an outline of algebra related to curvature tensor.
Most of the statements come without proofs, but everything can be proved by straightforward calculations (which are often tedious).

Most proofs can be found in \cite[Chapters 4+6]{do-carmo} and \cite[Chapter 3]{cheeger-ebin}.

\section{Definition}
Let $\vec x$, $\vec y$, $\vec v$, and $\vec w$ be vector fields on a Riemannian manifold $(M,g)$.
Recall that $\nabla$ denotes the Levi-Civita connection on $M$.

The \index{curvature tensor}\emph{Riemannian curvature tensor} $\Rm$ is defined by%
\footnote{Many authors (for example do Carmo) define it with opposite sign.
If you see notation $\mathop{\rm Rm}$ then most likely the sign is opposite.
This convention fits better with an earlier convention that sphere has positive Gauss curvature.}
\[\Rm(\vec x,\vec y)\vec v
=
\nabla_{\vec x}\nabla_{\vec y}\vec v
-\nabla_{\vec y}\nabla_{\vec x}\vec v
-\nabla_{[\vec x,\vec y]}\vec v.\]
It has valence 4 ---
it takes 3 vectors and spits another vector.
We do not need to specify covariance/contravariance type of the tensor since 
the metric tensor identifies tangent and cotangent bundles.

By the definition, one sees that the curvature tensor depends linearly on vector fields.
But it is indeed a tensor --- that is,  the vector $\Rm(\vec x,\vec y)\vec v$ depends only on the tangent vectors $\vec x$, $\vec y$, and $\vec v$ at the point.
The latter follows from the following identities: 
\begin{align*}
f\cdot \Rm(\vec x,\vec y)(\vec v)&=\Rm(f\cdot \vec x,\vec y)(\vec v)=
\\
&=\Rm(\vec x,f\cdot \vec y)(\vec v)=
 \\
&=\Rm(\vec x,\vec y)(f\cdot \vec v).
\end{align*}
for any vector fields $\vec x$, $\vec y$, $\vec v$ and a function $f$.
All of them can be proved by straightforward computations.

\section{Curvatuere transformation}

For given tangent vectors $\vec x$ and $\vec y$ at a point the linear map 
\[\Rm(\vec x,\vec y)\:\T\z\to\T\] is called \emph{curvature transformation}; some authors prefer to denote it by $\Rm_{\vec x,\vec y}$.
It has the following geometric meaning:

Let $\gamma$ be the contour of small parallelogram spanned by vectors $\vec x$ and $\vec y$ at a point $p$.
Let us denote by $\iota_\gamma\:\T_p\to\T_p$ the parallel transport along $\gamma$.
Denote by $a$ the area of parallelogram.
Then 
\[\iota_\gamma=\id+a\cdot \Rm(\vec x,\vec y) +o(a).\]
where $\id$ denotes the identity map on $\T_p$.

\begin{thm}{Exercise}
Suppose that parallel transport $\T_p\to \T_q$ in a Riemannian manifold $(M,g)$ does not depend on a path connecting $p$ to $q$.
Show that $(M,g)$ is \emph{flat}; that is, its curvature tensor vanish at all points.
\end{thm}


\section{Symmetries}

Set 
\[\hat\Rm(\vec x,\vec y,\vec v,\vec w)
\df
\langle \Rm(\vec x,\vec y)\vec v,\vec w\rangle.\]
Note that $\hat\Rm$ remembers everything about the curvature tenor $\Rm$ (assuming that metric tensor is known).
The $\hat \Rm$-form of $\Rm$ is more convenient to describe the symmetries of curvature tensor:
\[
\begin{aligned}
\hat\Rm(\vec x,\vec y,\vec v,\vec w)
=
-\hat\Rm(\vec y,\vec x,\vec v,\vec w)
=-\hat\Rm(\vec x,\vec y,\vec w,\vec v&),
\\
0=\hat\Rm(\vec x,\vec y,\vec v,\vec w)
+
\hat\Rm(\vec y,\vec v,\vec x,\vec w)
+
\hat\Rm(\vec v,\vec x)\vec y,\vec w&).
\end{aligned}
\eqlbl{eq:symR}
\]
The last identity is called the \emph{algebraic Bianchi identity} or \emph{first Bianchi identity} (and it was \emph{not} discovered by Bianchi).

These identities can be proved by straightforward computations.
Latter we will show that these symmetries provide a complete list; that is, given a tensor that satisfies these identiries, thare is a Riemannian manifold with such curvature tensor at some point.

The following equality follows from the main symmetries
\[
\hat \Rm(\vec x,\vec y,\vec v,\vec w)
=
\hat \Rm(\vec v,\vec w,\vec x,\vec y).
\eqlbl{eq:symR+}
\]

\section{Space of curvature tensors}\label{sec:A4}

Let us denote by $\A^4\T$ the space of all \emph{algebraic }curvature tensors on $\T$;
that is, $\A^4\T$ all valence 4 tenors with the symmetries \ref{eq:symR}.

Given a Euclidean space $E$, we denote by $\S^nE$ and $\Lambda^nE$ the space of symmetric and antisymmetric  tensors of valence $n$ over $E$.

Note that \ref{eq:symR+} and the first line in \ref{eq:symR} imply that 
\[\A^4\T\subset \S^2\Lambda^2\T.\]
In other words, a curvature tensor can be discribed as a symmetric bilinear form on the space of bivectors $\Lambda^2\T$;
or as the so-called \emph{curvature operator} --- a linear operator $\Rop\:\Lambda^2\T\z\to \Lambda^2\T$ defined by%
\footnote{If $\{\vec e_i\}$ is an orthonormal basis of $\T$,
then the scalar product on $\Lambda^2\T$ is defined by stating that it has an orthonormal basis $\{\vec e_i\wedge \vec e_j\}_{i>j}$.

Alternatively, the scalar product in $\Lambda^2\T$ can be also defined on simple bivectros $\vec x\wedge \vec y$ and $\vec v\wedge \vec w$ by stating that 
\[\langle\vec x\wedge \vec y,\vec v\wedge \vec w\rangle=\langle\vec x,\vec v\rangle\cdot \langle\vec y,\vec w\rangle-\langle\vec x,\vec w\rangle\cdot \langle\vec y,\vec v\rangle\]
and extended linearly to whole $\Lambda^2\T$.
} 
\[\langle \Rop(\vec x\wedge\vec y),\vec v\wedge\vec w\rangle=-\langle \Rm(\vec x,\vec y)\vec v,\vec w\rangle.\]
The symmetry \ref{eq:symR+} implies that $\Rop$ is self-adjoint; that is,
\[\langle\Rop\phi,\psi\rangle=\langle\phi,\Rop\psi\rangle\]
for any bivectors $\phi,\psi\in\Lambda^2\T$.

The algebraic Bianchi identity implies that complete antisymmetrization of $\hat \Rm$ vanish. 
More precisely, if $\alpha\:\S^2\Lambda^2\T\to\Lambda^4\T$ denotes the complete anysymmetrization then space of curvature tensors is the kernel of $\alpha$.%
\footnote{As far as I see, the following property is absolutely useless, but it is funny.
Given a curvature tenor $\Rm$ consider the tensor 
\[\Ya(\vec x,\vec v,\vec y,\vec w)\df\langle \Rm(\vec x,\vec y)\vec v,\vec w\rangle+\Rm(\vec v,\vec y)\vec x,\vec w\rangle.\]
Then the linear transformation $\Rm\to \Ya$ describes an isomorphism 
\[\S^2\Lambda^2\T\ominus\Lambda^4\T
\longleftrightarrow
\S^2\S^2\T\ominus\S^4\T.\]
That is, curvature tensor can described as quadratic forms on quadratic forms that lie in the kernal of complete symmetrization $\sigma\:\S^2\S^2\T\z\to \S^4\T$;
in other words $\Ya$ satisfies the following symmetries:
\begin{align*}
\Ya(\vec x,\vec v,\vec y,\vec w)
=
\Ya(\vec v,\vec x,\vec y,\vec w)
&=
\Ya(\vec x,\vec v,\vec w,\vec y)=\Ya(\vec y,\vec w,\vec x,\vec v),
\\
0&=\Ya(\vec x,\vec x,\vec x,\vec x).
\end{align*}
}
The latter can be written as 
\[\A^4\T=\S^2\Lambda^2\T\ominus\Lambda^4\T
\quad\text{or}\quad
\A^4\T=\S^2\Lambda^2\T\cap(\Lambda^4\T)^\perp,\]
where $L^\perp$ stands for the orthogonal complement $L$ in the Euclidean metric on $\S^2\Lambda^2\T$ induced from $\T$.

If $n=\dim\T$ then the dimension of the space of curvature tensors over $\T$ can be easily calculated:
\[\dim (\S^2\Lambda^2\T)- \dim(\Lambda^4\T)=
\binom{\binom n2+1}2-\binom n4=\frac{n^2\cdot(n^2-1)}{12}.\]

\begin{thm}{Exercise}\label{ex:all-tensors}
Suppose $\Rop_\phi:\Lambda^2\T\to\Lambda^2\T$ is an othogonal projection to a 1-dimnsional subspace spanned by a bivector $\phi\in \Lambda^2\T$.
Show that $\Rop$ is a curvature operator if and only if $\phi$ is a simple bivector;
that is, if $\phi=\vec x\wedge\vec y$ for some $\vec x,\vec y\in\T$.

Show that in this case $\Rop$ is the curvature operator of $\SSS^2\times \RR^{n-2}$.

Use the results in Section \ref{sec:submanifolds} to show that any algebraic curvature tensor can appear as a curvature tensor of a Riemannian manifold.
\end{thm}



\section{Sectional curvature}

Let $p$ be a point in a Riemannian manifold $(M,g)$.
Choose a plane $\sigma$ in the tangent space $\T_p$. 
Consider a surface $\Sigma$ in $M$ sweeped by short geodesics from $p$ in the directions on $\sigma$.
The Gauss curvature of $\Sigma$ at $p$ is called \emph{sectional curvature} and denoted by $\sec\sigma$;
the plane $\sigma$ is called \emph{sectional direction}.

If the sectional direction $\sigma$ is spanned by vectors $\vec x$ and $\vec y$, then
\begin{align*}
\sec\sigma&=\frac{\langle \Rm(\vec x,\vec y)\vec y,\vec x\rangle}{|\vec x|^2\cdot|\vec y|^2-\langle\vec x,\vec y\rangle^2}=
\\
&=\frac{K(\vec x,\vec y)}{|\vec x\wedge\vec y|^2};
\end{align*}
in the last expression we use shortcut
\[K(\vec x,\vec y)=\langle \Rm(\vec x,\vec y)\vec y,\vec x\rangle=-\langle \Rm(\vec x,\vec y)\vec x,\vec y\rangle=\langle\Rop(\vec x\wedge\vec y),\vec x\wedge\vec y\rangle;\]
note that $K$ is quadratic in both arguments.

The formula above implies that sectional curvature can be recovered from curvature tensor.

The curvature tensor can be expressed using sectional curvature as well.
Indeed, once we know curvature in any sectional direction, we can use the above formula to find $K(\vec x,\vec y)$ for any tangent vectors $\vec x,\vec y\in\T_p$.
After that one could apply the following  formula:
\begin{align*}
6\cdot \Rm(\vec x,\vec y)\vec v,\vec w\rangle&=
[K(\vec x+\vec v,\vec y+\vec w)
+K(\vec x,\vec w)
+K(\vec y,\vec v)-\\
-K(\vec x&+\vec v,\vec y)
-K(\vec x+\vec v,\vec w)
-K(\vec x,\vec y+\vec w)
-K(\vec v,\vec y+\vec w)
]-
\\
&-[K(\vec x+\vec w,\vec y+\vec v)
+K(\vec y,\vec w)
+K(\vec x,\vec v)
\\
-K(\vec x&+\vec w,\vec y)
-K(\vec x+\vec w,\vec v)
-K(\vec x,\vec y+\vec v)
-K(\vec w,\vec y+\vec v)].
\end{align*}
The formula is scary, but it is very similar to recovery of a symmetric bilinear form its quadratic form.

\begin{thm}{Exercise}\label{ex:sec/Rop}
Let $\dim\T=3$.
Suppose that a a curvature tenor $\Rm\in\A^4\T$ has positive sectional curvature in all directions.
Show that the corresponding curvature operator is positive;
that is $\langle\Rop\phi,\phi\rangle>0$ is for any $\phi\in\Lambda^2\T$. 

Show that if $\dim\T=4$, then analogous statement does not hold.
\end{thm}


\section{Ricci decomposition}

Let ${\vec e_i}$ be an orthonormal basis at a point $p$ of Riemannian manifold.
The so called \emph{Ricci curvature tensor} is a linerar transformation, $\Ric\:\T_p\to\T_p$ defined as
\[\langle\Ric\vec x,\vec x\rangle=\sum_i K(\vec x,\vec e_i).\]
Further, the \emph{scalar curvature} $\Sc$ is defined by 
\[\Sc=\sum_{i,j} K(\vec e_i,\vec e_j)=\sum_{j}\langle\Ric \vec e_j,\vec e_j\rangle=2\cdot\trace \Rop.\]

If $|\vec x|=1$ then the value $\langle\Ric(\vec x),\vec x\rangle$ is called \emph{Ricci curvature} in the direction $\vec x$.
For example $n$-dimensional unit sphere has sectional curvature 1, Ricci curvature $n-1$ in all directions and scalar curvature $n\cdot(n-1)$.

The action of orthogonal group $O(\T)$ can be extended to $\A^4\T$.
Evidently the kernels of $\Rm\to \Ric$ and $\Rm\to\Sc$ and their orthogonal complements are invariant with respect to this action.
In other words, Ricci tenor $\Ric$ and scalar curvature $\Sc$ do not depend on the choice of the orthonormal basis ${\vec e_i}$.

It turns out that these are the only subspaces of $\A^4\T$ that invariant with respect to $O(\T)$.
In other words, the decomposition 
\[\A^4\T=U\oplus V \oplus W\]
with the subspaces $U$, $V$, and $W$ described below is the maximal $O(\T)$-invarint decomposition of $\A^4\T$.
\begin{itemize}
\item $W$ is the kernel of $\Rm\to \Ric$; the orthogonal projection to $W$ of the curvature tensor is called its \emph{Weil tensor}.
\item $V$ is the intersection of the kernel $\Rm\to\Sc$ and the orthogonal complement of $W$. The projection to $V$ is completely determined by the \emph{traceless Ricci tenor} $\Ric_0=\Ric-\tfrac1n\cdot \Sc\cdot g$.
\item $U$ is a one-dimensional subset of curvature tenors with constant sectional curvatures; so the curvature operator is proportional to the identity.
\end{itemize}

 



\section{Curvature bounds}

The following theorem is a classical result in Riemannian geometry with long history.

\begin{thm}{Quarter-pinched sphere theorem}
Suppose that $(M,g)$ is a simply connected Riemannian manifold with sectional curvature strictly between $1$ and $4$ at each point.
Then $M$ is diffeomorphic to a sphere.
\end{thm}

It gives an example of the so called \emph{local-to-global theorems}.
Its assumption is a local property that is given by a curvature bounds at each point.
The conclusion says something about global structure of the manifold; in this example it says something about its topological type.

Riemannian geometry has other types of theorems\footnote{for example \emph{rigidity theorems} say that a Riemannian manifold with given property must be isometric to some known manifold (typically a round sphere).}, but nothing else makes Riemannian geometer nearly as happy as the local-to-global results.

Typically the local condition is given by restriction on curvature tensor of Riemannian manifold;
that is, we specify a subset $\Omega\subset \A^4\T$ and assume that curvature tensor of a Riemannian manifold belongs to $\Omega$ at each point.
Most of the time the set $\Omega\subset \A^4\T$ is open or at least has nonempty interior.\footnote{There are exceptions, for example Einstein manifolds defined by an equation on curvature tensor.
But this subject lies on a half way from differential geometry to partial differential equations.}

It is reasonable to assume that $\Omega$ is convex (or at least connected).
If $\dim\T=2$, then $\dim\A^4\T=1$;
in this case we do not have much choice --- we might consider curvature bounded above, or below, or from both sides.
That what we did for surfaces.

Starting from dimension 3, things getting more complicated --- it is reasonable to assume in addition that $\Omega$ is invariant with respect to rotations of the tangent space;
in other words $\Omega$ has to be invariant with respect to action of the orthonormal group $O(\T)$ of the space $\T$ extended to $\A^4\T$.
But still, we have huge family of curvature conditions.

The Ricci flow technique deals with few families of curvature bounds in the proofs.
Couple of dozens of curvature bounds made it to a formulation a meaningful theorem.
The champions seem to be upper, lower, and bilateral bounds on sectional curvature, lower bounds on Ricci curvature, and lower bounds on scalar curvature.%
\footnote{As far as I see, the following property is absolutely useless, but it is funny.
The cone of curvature tensors with nonnegative (or nonpositive) sectional curvature is a maximal $\GL(\T)$-invariant subsets of $A^4(\T)$.}
 
\section{Submanifolds}\label{sec:submanifolds}

Suppose $M$ is a submanifold in a Riemannian manifold $(\bar M,g)$.
Recall that restricting $g$ to the tangent bundle over $M$ produce a Riemannian metric on $M$.
Let us denote by $\bar\nabla $ and $\nabla$ the Levi-Civita connction on $\bar M$ and $M$ respectively.

Recall that second fundamental form $S$ of $M$ is defined by
\[S(\vec x,\vec y)=\nabla_{\vec x}\vec y-\bar\nabla_{\vec x}\vec y;\]
it takes two tangent vectors on $M$ and spits a normal vector;
so is a tensor in $\S^2\T M\otimes \N M$ (here we identify tangent/cotangent normal/conormal bundles as usual).

Again straightforward computations show that that $S$ is indeed a tensor --- one has to show that 
\[f\cdot S(\vec x,\vec y)=S(f\cdot \vec x,\vec y)=S(\vec x,f\cdot \vec y)\]
for any tangent vector fields $\vec x$, $\vec y$, and a smooth function $f$ on $M$.

The following formula gives a relation between curvature tensors $\Rm$ and $\bar \Rm$ of $M$ and $\bar M$ respectively:
\begin{align*}
\langle \Rm(\vec x,\vec y)\vec v,\vec w\rangle
&=
\langle \bar \Rm(\vec x,\vec y),\vec v,\vec w\rangle
+
\\
&+
\langle S(\vec x,\vec w),S(\vec y,\vec v)\rangle
-
\langle S(\vec x,\vec v),S(\vec y,\vec w)\rangle
.\end{align*}
In particular, using the shortcut $K(\vec x,\vec y)=\langle \Rm(\vec x,\vec y)\vec y,\vec x\rangle$, we can write
\[K(\vec x,\vec y)=\bar K(\vec x,\vec y)+\langle S(\vec x,\vec x),S(\vec y,\vec y)\rangle
-
|S(\vec x,\vec y)|^2.\]

This is a generalization of the formula for Gauss curvature of surface in the Euclidean space:
\[K=\ell\cdot n-m^2,\]
where $\ell$, $m$, and $n$ are components of Hessian matrix at $p$ in an othonormal basis.
Indeed, if $\vec x,\vec y$ is an orthonormal frame at a point in a surface, then $K=K(\vec x,\vec y)$ is its Gauss curvature, and
\begin{align*}
S(\vec x,\vec x)&=\ell\cdot\Norm,
&
S(\vec x,\vec y)&=m\cdot\Norm,
&
S(\vec y,\vec y)&=n\cdot\Norm,
\end{align*}
where $\Norm$ is a unit normal vector at $p$.

\begin{thm}{Exercise}\label{ex:4-hypersurface}
Show that there is a 4-dimensional Riemannian manfifold $(M,g)$ such that no neighborhood of a point $p$ in $M$ admits a smooth length-preserving embedding in $\RR^5$.
\end{thm}

\parit{Hint:} Count dimensions of second fundamental form and curvature tensor at $p$ and apply the formula.

\begin{thm}{Exercise}\label{ex:codim=2}
Let $M$ be a smooth submanifold of a Euclidean space.
Assume $\codim M=2$.
Show that if sectional curvature of $M$ is positive, then so is its curvature operator.

An analogues statement for submanifolds of codimension 3 does not hold --- try to guess why.
\end{thm}

\parit{Hint:} 
Show first that 
$\langle S(\vec x,\vec x),S(\vec y,\vec y)\rangle> 0$
for any two nonvanishing tangent vectors $\vec x$, $\vec y$ at any point $p\in M$. 
Furhter, show and use that if $\codim M=2$, then the normal space admits an orthonormal basic $\vec u_1,\vec u_2$ such that both quadratic forms
$s_i(\vec x,\vec y)\df\langle S(\vec x,\vec y),\vec u_i\rangle$
are positive definite. 


\section{Submersions}

Let $s\:\bar M\to M$ be a submersion between smooth manifolds.
Suppose that manifolds $\bar M$ and $M$ are equipped with Riemannian metrics.

Suppose that $s(\bar p)=p$.

The kernel of the differential $ds\:\T_{\bar p} \bar M\to \T_p M$ will be called \emph{vertical subspace};
it will be denoted by $\V_{\bar p}$.
The orthogonal complement of $\V_{\bar p}$ in $\T_{\bar p}$ will be called \emph{horisontal subspace} at ${\bar p}$;
it will be denoted by $\Hor_{\bar p}$.

The submersion $s$ is called \emph{Riemannian} if the restriction of $ds$ to $\Hor$ is an isometry.

Note that $\T_{\bar p}=\Hor_{\bar p}\oplus \V_{\bar p}$.
In particular, any tangent vector $\vec x\z\in\T\bar M$ can be spitted into its vertical and horizontal part denoted by $\vec x^\Hor$ and $\vec x^\V$, so 
\[\vec x=\vec x^{\Hor}+\vec x^\V.\]

By the definition of Riemannian submersion,
for any vector $\vec x\in\T_p$ there is a unique vector $\bar{\vec x}\in\Hor_{\bar p}$ such that 
$\vec x=ds(\bar{\vec x})$.

The curvature of $ M$ can be found using the so called \emph{O'Neil formula}:
\[K(\vec x,\vec y)=\bar K(\bar{\vec x},\bar{\vec y})+\tfrac34\cdot|[\bar{\vec x},\bar{\vec y}]^\V|^2.\]

The notion of Riemannian submersion is a dual to submanifold ($\approx$ lenght-preserving immersion).
The tensor $A$ defined by
\[A(\bar{\vec x},\bar{ \vec y})=[\bar{\vec x},\bar{ \vec y}]^\V\]
is indeed a tensor in $\Lambda^2\Hor\otimes\V$ (it is proved the usual way).
This tensor plays a role similar to the second fundamental form of a submanifold.


\begin{thm}{Exercise}\label{ex:4-form}
Let $\bar M\to M$ be a Riemannian submersion.
Suppose that $\bar M$ has nonnegative curvature operator at all points.
Show that at any point of $M$ there is a 4-vector $\eta$  such that 
\[\langle\Rop \phi,\phi\rangle+\langle\eta,\phi\wedge\phi\rangle\ge 0\]
for any 2-vector $\phi$.
\end{thm}

\parit{Hint:} Use Section~\ref{sec:A4}.

\section{Lie groups}

Suppose $G$ is a Lie group with biinvariant metric.
By straightforward comutations, one gets the following identities for left-invariant vector fields on $G$:
\begin{align*}
\nabla_{\vec x}\vec y&=\tfrac12\cdot[\vec x,\vec y]
\\
\langle\Rm(\vec x,\vec y)\vec v,\vec w\rangle&=\tfrac14\cdot(\langle [\vec x,\vec w],[\vec y,\vec v]\rangle-\langle[\vec x,\vec v],[\vec y,\vec w]\rangle)
\\
K(\vec x,\vec y)&=\tfrac14\cdot|[\vec x,\vec y]|^2
\end{align*}

Note that the first identity implies that $\nabla_{\vec x}\vec x=0$.
Therefore homomorphisms $\RR\to G$ are geodesics.


\section{Cheeger's trick}

Suppose a Lie group $G$ acts isometrically on a Riemannian manifold $M=(M,g)$.
Suppose that a $G$ admits a bi-invarinat metric (this is alway the case if $G$ is compact).

Consider the diagonal action of $G$ on the product $G\times M$;
that is, $a\cdot (b,x)\df(a\cdot b, a\cdot x)$ for any $a,b\in G$ and $x\in M$.
Note that this action is isometric and free.
Therefore the quotient map $(\lambda\cdot G)\times M\to (M,g_\lambda)$ is a Riemannian sumersion for some metric $g_\lambda$.
This way we obtain a one parameter family $g_\lambda$ of Riemannian metrics on $M$.
Note that $g_\lambda\to g$ as $\lambda\to \infty$.

This procedure is called \emph{Cheeger's trick} and the obtained family is called \emph{Cheeger's deformation}.
(Jeff Cheeger found number of applications of this trick, but it was invented earlier.)

By O'Nail formula, if $(M,g)$ had a nonnegative (respectively positive) sectional curvature, then so does $(M,g_\lambda)$ for any $\lambda>0$.

\section{Berger spheres}

Applying the Cheegers to the isometric action of $\SSS^1$ on $\SSS^3$ by complex multiplication one gets \emph{Berger spheres} --- a family of metrics $g_\lambda$ on $\SSS^3$ with positive sectional curvature.

The Berger spheres can be also described as certain left-invariant metrics on on the Lie group $\SSS^3=\Spin(4)$.
There are explicit formulas for connection and curvature in such metrics.

It is an important example in Riemannian geometry.
(It could be compared to the Cantor set in analysis.)

\begin{thm}{Exercise}
Suppose $(\SSS^3,g_\lambda)$ denotes Berger spheres with parameter $\lambda$.
Set $\ell=\lambda/\sqrt{1+\lambda^2}$.

Show that

\begin{subthm}{}
$(\SSS^3,g_\lambda)$ are foliated by closed geodesics of length $2\cdot\pi\cdot\ell$
and 
\[\vol(\SSS^3,g_\lambda)=\ell\cdot\vol\SSS^3.\]
\end{subthm}

\begin{subthm}{}
Show that the curvature tensor of $(\SSS^3,g_\lambda)$ converges to the curvature tensor of $\SSS^2_{\frac12}\times\RR$ as $\lambda\to0$, where $\SSS^2_{\frac12}$ denotes a round sphere of radius $\tfrac 12$.
(One has to find right sense for term ``converge'' here.)
\end{subthm}

\begin{subthm}{}
Try to visualize what happens with  $(\SSS^3,g_\lambda)$ as $\lambda\to0$.
\end{subthm}


 \end{thm}




